% ==============================
% Chapter 1 – Introduction
% ==============================

\chapter{Introduction}

% --------------------------------
% 1.1 Introduction
% --------------------------------

\section{Introduction}

Writing software for quantum computers today resembles writing machine code for mainframes in the late 1950s. While theoretical quantum computing promises significant speedups for prime factorization, quantum chemistry, and combinatorics \cite{nielsen2010}, the practical programming landscape remains fragmented across incompatible physical backends. A classical programmer compiles high-level code down to a standardized Instruction Set Architecture (ISA) like x86 or ARM, protected by operating system abstractions. Quantum programmers, in contrast, must still program directly against the physical constraints of individual dilution refrigerators and ion traps.

This fragmentation stems directly from the physics of quantum hardware. During the current Noisy Intermediate-Scale Quantum (NISQ) era \cite{preskill2018}, hardware vendors pursue distinct physical realizations:
\begin{itemize}[leftmargin=*]
    \item \textbf{Superconducting Transmons:} IBM and Google fabricate lithographic circuits on silicon, where fixed or tunable microwave pulses drive Josephson junctions. Qubits are laid out in planar planar graphs, such as IBM's heavy-hex lattice or Google's Sycamore grid. Two-qubit gates can only fire between physically connected neighbours.
    \item \textbf{Trapped Ions:} Vendors like IonQ trap ionized ytterbium or barium atoms in linear radio-frequency Paul traps. Laser beams steer interactions between arbitrary ion pairs, yielding near-all-to-all connectivity, but gate switching times are orders of magnitude slower than solid-state devices.
    \item \textbf{Neutral Atoms and Photonics:} Companies such as QuEra and Xanadu use optical tweezers or spatial beam splitters, exposing entirely different interaction models and native basis operations.
\end{itemize}

Because these systems diverge in native gate sets, coupling topologies, and coherence times, quantum algorithms cannot simply run across backends without manual rewiring. A circuit tuned for IBM's native basis $\{RZ, \sqrt{X}, X, CX\}$ and heavy-hex layout fails completely on Google's $\{XY, CZ\}$ Sycamore processor unless someone manually decomposes every gate, inserts routing SWAPs, and accounts for hardware-specific constraints.

Current industry frameworks only partially alleviate this friction. Software development kits like Qiskit \cite{qiskit2023}, Cirq \cite{cirq2020}, and PyQuil offer compilation tools, but they tend to pull heavy enterprise dependencies and remain closely tied to their sponsors' hardware and cloud services. Installing Qiskit requires pulling hundreds of megabytes of compiled binaries, Rust-based extensions, and scientific libraries. For undergraduate students, educators, and independent researchers, setting up these ecosystems just to test a 5-qubit algorithm or inspect compilation steps creates an unnecessarily high barrier to entry.

These friction points motivated our Final Year Project: the \textbf{Quantum Virtual Machine (QVM)}. QVM is a lightweight, hardware-agnostic compilation and execution runtime designed around the ``Write Once, Run Anywhere'' (WORA) philosophy. Rather than writing code for one proprietary cloud service, a developer writes their algorithm in standard OpenQASM 2.0 \cite{cross2017qasm}, OpenQASM 3.0 \cite{cross2022qasm3}, or JSON. QVM parses the circuit, converts it into an internal hardware-independent Intermediate Representation (IR), decomposes non-native gates, and maps qubits onto arbitrary coupling graphs using heuristic routing algorithms.

To verify programs without cloud queues or API fees, QVM includes two complementary execution backends. An exact statevector simulator handles circuits up to 12 qubits with microsecond-level gate updates using vectorized NumPy arrays \cite{numpy2020}. When circuits expand beyond 12 qubits, a Matrix Product State (MPS) tensor network simulator \cite{vidal2003mps} steps in, using Singular Value Decomposition (SVD) truncation to simulate 20+ qubits as long as circuit entanglement remains bounded. By pairing compilation transparency with dual simulation engines, QVM gives students and engineers a clear view into every step of the quantum execution pipeline.

% --------------------------------
% 1.2 Vision Statement
% --------------------------------

\section{Vision Statement}

Our vision for the Quantum Virtual Machine is straightforward: isolate quantum algorithmic logic from the physical layout of quantum hardware.

When a student or researcher designs an algorithm, their primary focus should be quantum logic: superposition, interference, and phase kickback. They should not need to spend days debugging topological route failures or rewriting gate decompositions every time a hardware vendor updates its chip revision. 

QVM aims to serve as an accessible educational and research platform that bridges three core audiences:
\begin{itemize}[leftmargin=*]
    \item \textbf{Computer Science Students:} Learners need to see under the hood. Commercial SDKs often hide routing and decomposition steps behind complex compiler passes. QVM exposes each stage clearly: parsing, AST construction, intermediate representation, SWAP routing, and statevector evolution.
    \item \textbf{Algorithm Researchers:} Researchers comparing routing strategies need a quick, reproducible testbed. QVM allows benchmarking Greedy BFS versus SABRE routing on identical circuits without cloud latency or hardware fees.
    \item \textbf{Software Integrators:} Developers building quantum frontends or toolchains need a modular backend that can be installed in seconds via \texttt{pip}, run locally without cloud credentials, and integrated through a RESTful API.
\end{itemize}

By decoupling algorithm representation from hardware realization, QVM treats target chips as pluggable constraint graphs rather than proprietary silos.

% --------------------------------
% 1.3 Related System Analysis
% --------------------------------

\section{Related System Analysis}

To place QVM in context, we examined four prominent quantum programming frameworks. Each tool brings notable strengths, but each also introduces distinct trade-offs in complexity, dependency footprint, or vendor focus.

\textbf{Qiskit (IBM Quantum)} \cite{qiskit2023}: Qiskit is the dominant framework in the quantum open-source ecosystem. It provides extensive libraries for quantum chemistry, optimization, machine learning, and hardware pulse calibration. Its transpiler features sophisticated optimization pipelines. However, this breadth has made Qiskit heavy. A full installation routinely exceeds 500~MB and pulls deep dependency trees. For beginners, the abstraction layers can feel opaque. It is not always obvious what physical transformations occurred between a raw circuit and the compiled output. Additionally, while third-party plugins exist, Qiskit's primary development priority naturally tracks IBM hardware.

\textbf{Cirq (Google Quantum AI)} \cite{cirq2020}: Cirq was built from the ground up for near-term experimental quantum algorithms on Google's Sycamore architecture. It prioritizes exact control over gate scheduling, moment structure, and spatial qubit placement. This low-level precision appeals to physicists running pulse experiments, but it requires programmers to handle manual layout and scheduling details that high-level compilers usually automate. Cirq's circuit construction syntax is verbose, and porting circuits to topologies outside Google's planar grid requires considerable glue code.

\textbf{ProjectQ (ETH Zurich)} \cite{haner2018}: ProjectQ was an early pioneer in modular quantum compilation. It explicitly separated high-level circuit syntax from backend compilation engines, demonstrating the power of a compiler-pipeline model. However, active open-source maintenance of ProjectQ has slowed significantly. It lacks native support for OpenQASM 3.0 syntax, classical conditional branching, and modern hybrid control flow.

\textbf{PyQuil (Rigetti Computing)}: PyQuil exposes Rigetti's Quil language, emphasizing hybrid quantum-classical computing through its Quantum Abstract Machine (QAM) interface. PyQuil excels at parametric loops for variational algorithms, but it remains centered around Rigetti's Quil ecosystem and Forest cloud service. Learning Quil introduces another vendor-specific syntax rather than standardizing on portable formats like OpenQASM.

Table \ref{tab:related_systems} summarizes the comparative strengths and weaknesses of these systems alongside the specific design improvements implemented in QVM.

\vspace{0.5cm}

\begin{table}[h!]
\centering
\caption{Comparative Analysis of Quantum Computing Frameworks}
\label{tab:related_systems}
\renewcommand{\arraystretch}{1.3}

\begin{tabularx}{\textwidth}{|p{3cm}|X|X|}
\hline
\textbf{Framework} & \textbf{Observed Limitations} & \textbf{QVM Architectural Approach} \\
\hline
Qiskit (IBM) &
\begin{itemize}[leftmargin=*, noitemsep, topsep=0pt]
    \item Installation footprint exceeds 500~MB
    \item Deep dependency tree
    \item High conceptual complexity for beginners
    \item Strong alignment with IBM hardware backends
\end{itemize}
&
\begin{itemize}[leftmargin=*, noitemsep, topsep=0pt]
    \item Lightweight footprint ($<50$~MB)
    \item Minimal core dependencies (NumPy, Lark)
    \item Clear, stepwise compilation inspection
    \item Pluggable, vendor-neutral hardware topologies
\end{itemize}
\\
\hline
Cirq (Google) &
\begin{itemize}[leftmargin=*, noitemsep, topsep=0pt]
    \item Requires explicit manual moment scheduling
    \item Circuit syntax is verbose for simple routines
    \item Tailored primarily to Google Sycamore grids
\end{itemize}
&
\begin{itemize}[leftmargin=*, noitemsep, topsep=0pt]
    \item Automated gate decomposition and routing
    \item Standard OpenQASM 2.0 and 3.0 parsing
    \item Generic graph representation for any topology
\end{itemize}
\\
\hline
ProjectQ (ETH) &
\begin{itemize}[leftmargin=*, noitemsep, topsep=0pt]
    \item Stagnant development activity
    \item No native OpenQASM 3.0 grammar support
    \item Lacks modern classical control flow
\end{itemize}
&
\begin{itemize}[leftmargin=*, noitemsep, topsep=0pt]
    \item Active maintenance with Python 3.10+
    \item AST-based OpenQASM 3.0 parser via Lark
    \item Full classical registers, conditionals, and loops
\end{itemize}
\\
\hline
PyQuil (Rigetti) &
\begin{itemize}[leftmargin=*, noitemsep, topsep=0pt]
    \item Dependent on Rigetti Quil syntax
    \item Tight coupling with Forest cloud platform
    \item Limited out-of-the-box cross-platform tools
\end{itemize}
&
\begin{itemize}[leftmargin=*, noitemsep, topsep=0pt]
    \item Native OpenQASM input/output compatibility
    \item Fully self-contained local execution
    \item Zero mandatory cloud accounts or API keys
\end{itemize}
\\
\hline
\end{tabularx}
\end{table}

% --------------------------------
% 1.4 Project Deliverables
% --------------------------------

\section{Project Deliverables}

To realize the goals of the project, we engineered a set of modular software components, verification suites, and technical documentation:

\begin{enumerate}[leftmargin=*]
    \item \textbf{Core QVM Compilation Engine:} The central processing pipeline spanning six modular layers: Input, Parsing, Intermediate Representation, Gate Decomposition, Transpilation, and Simulation. It coordinates the journey from raw circuit text to final quantum probability amplitudes.
    \item \textbf{Grammar-Based Parsers:} A formal LALR(1) parser generated via Lark \cite{lark2023} for OpenQASM 3.0, paired with an OpenQASM 2.0 regex parser and a JSON format loader. The parser validates syntax, detects malformed declarations, and builds structured Abstract Syntax Trees (ASTs).
    \item \textbf{Hardware Transpiler:} A topology-aware routing system implementing both Greedy BFS and SABRE \cite{li2019sabre} heuristics. It automatically detects non-adjacent two-qubit operations, computes optimal routing paths, and inserts SWAP gates with optional post-execution mapping restoration.
    \item \textbf{Dual Simulation Backends:} (1) An exact statevector simulator optimized with NumPy array slicing for circuits up to 12 qubits, supporting mid-circuit measurements and classical bit registers. (2) An MPS tensor network simulator scaling to 20+ qubits for low-entanglement states using SVD singular value truncation.
    \item \textbf{Kraus Noise Modeling Module:} An open-system noise engine simulating real-world NISQ imperfections, including depolarizing channels, amplitude damping ($T_1$), phase damping ($T_2$), and asymmetric readout confusion matrices.
    \item \textbf{Command-Line Interface (CLI):} A terminal application with arguments for file loading, backend selection, routing strategy configuration, shot-based sampling, and Matplotlib plotting.
    \item \textbf{FastAPI Execution Service:} A RESTful HTTP backend providing asynchronous endpoints (\texttt{POST /run}) for circuit simulation, coupled with Supabase cloud database persistence for saving circuits and retrieving execution history.
    \item \textbf{Web Dashboard:} A responsive browser application built with Next.js 16, React 19, and Tailwind CSS. It features an in-browser circuit editor, preset algorithm loaders, live SVG circuit schematics, and interactive probability charts.
    \item \textbf{Automated Test Suite:} A test suite of 215 automated pytest tests verifying simulator numerical stability, parser error handling, routing invariants, and canonical quantum algorithms.
    \item \textbf{Verified Algorithm Implementations:} End-to-end demonstrations of Bell state preparation, GHZ-3 state generation, Bernstein-Vazirani parity search, Grover's search algorithm, and Variational Quantum Eigensolver (VQE) optimization.
    \item \textbf{Documentation Artifacts:} Complete software specification and design documentation, user guides, API references, and compilation instructions.
\end{enumerate}

% --------------------------------
% 1.5 System Limitations / Constraints
% --------------------------------

\section{System Limitations / Constraints}

Developing a practical quantum simulator on classical consumer hardware requires navigating strict mathematical and physical boundaries:

\textbf{Exponential Memory Scaling:} The statevector of an $n$-qubit quantum system exists in a Hilbert space of dimension $2^n$. Because each complex amplitude requires 16 bytes (two 64-bit floating-point numbers), memory consumption doubles with every single added qubit. A 10-qubit circuit needs only 16~KB of memory; a 20-qubit circuit needs 16~MB; a 30-qubit circuit demands 16~GB; and a 40-qubit circuit would require 16 terabytes. Consequently, our exact statevector simulator is practically bounded to roughly 12 to 14 qubits on standard developer laptops with 8--16~GB RAM.

\textbf{Entanglement Limits in MPS Simulation:} The Matrix Product State simulator bypasses exponential memory bounds by approximating states as a one-dimensional chain of low-rank tensors. For weakly entangled circuits, this scales efficiently to 20+ qubits. However, when quantum circuits produce high, global entanglement (such as deep random circuit sampling or multi-qubit Quantum Fourier Transforms), the required bond dimension $\chi$ explodes exponentially. Once $\chi$ grows too large, the computational cost of tensor contractions approaches or exceeds dense statevector simulation.

\textbf{Heuristic Transpilation Optimality:} Optimal qubit placement and routing on restricted graphs is an NP-hard combinatorial problem. Both Greedy BFS and SABRE are heuristic algorithms. While SABRE achieves substantial depth reductions through lookahead scoring, it does not guarantee the mathematically minimum number of SWAP gates.

\textbf{Noiseless Baseline Assumption:} Unless the user explicitly activates the Kraus noise module, the simulator computes idealized unitary evolution. Physical quantum hardware exhibits time-dependent drift, crosstalk, and thermal fluctuations that idealized simulation does not capture by default.

\textbf{Classical Emulation Overhead:} Transpilation and classical simulation carry runtime overhead. Transpiling deep circuits with extensive lookahead windows requires evaluating graph distances across hundreds of gates, which is computationally heavier than running a trivial script. QVM is intended for algorithm development, educational exploration, and verification, rather than pretending to out-compute real supercomputers.

% --------------------------------
% 1.6 Tools and Technologies
% --------------------------------

\section{Tools and Technologies}

Every tool in the QVM stack was chosen to maintain a small installation footprint, reliable execution, and clear code organization:

\begin{itemize}[leftmargin=*]
    \item \textbf{Python 3.10+:} Selected as the core language for its dominance in the quantum community, clean syntax, and mature scientific libraries. Version 3.10+ provides structural pattern matching and enhanced type hints, improving codebase clarity.
    \item \textbf{NumPy (1.24+):} Provides core numerical computation. Instead of naive nested loops, statevector updates use NumPy's C-accelerated array slicing and axis permutations, keeping single-qubit gate operations under a few microseconds.
    \item \textbf{Lark (1.1+):} A pure-Python parsing toolkit used to compile OpenQASM 3.0 grammars into ASTs. Lark was chosen over ANTLR and PLY because it requires zero external runtimes, generates zero fragile C extensions, and delivers informative syntax error messages.
    \item \textbf{FastAPI (0.100+) and Uvicorn:} Powers the asynchronous REST API. FastAPI uses Pydantic to validate circuit payloads automatically, generates live OpenAPI interactive documentation, and handles requests without blocking simulation worker tasks.
    \item \textbf{Matplotlib (3.7+):} Generates static circuit diagrams and measurement histograms for reports and publications.
    \item \textbf{Next.js 16 and React 19:} Forms the web frontend, combining server-rendered static layout with client-side reactive editors, KaTeX mathematical typesetting, and responsive SVG circuit rendering.
    \item \textbf{pytest (7.4+):} Drives test automation with fixtures, parametrized test runs, and floating-point tolerance assertions via \texttt{numpy.allclose()}.
    \item \textbf{Git and GitHub:} Handles distributed version control, issue tracking, and release management.
\end{itemize}

Table \ref{tab:tools_tech} lists the versions and primary roles of each technology used in the project.

\vspace{0.4cm}

\begin{table}[h!]
\centering
\caption{System Technology Stack}
\label{tab:tools_tech}
\renewcommand{\arraystretch}{1.3}
\begin{tabularx}{\textwidth}{|p{3.5cm}|p{2.5cm}|X|}
\hline
\textbf{Tool / Technology} & \textbf{Version} & \textbf{Role in QVM} \\
\hline
Python & 3.10+ & Primary implementation language \\
\hline
NumPy & 1.24+ & High-performance linear algebra and array indexing \\
\hline
Lark & 1.1+ & Formal EBNF grammar parsing for OpenQASM 3.0 \\
\hline
FastAPI & 0.100+ & Asynchronous REST API service \\
\hline
Matplotlib & 3.7+ & Vector circuit diagram generation and probability plotting \\
\hline
Next.js / React & 16 / 19 & Interactive web user interface and SVG schematic viewer \\
\hline
pytest & 7.4+ & Test automation and numerical tolerance verification \\
\hline
Git / GitHub & 2.40+ & Version control and collaborative code maintenance \\
\hline
\end{tabularx}
\end{table}

% --------------------------------
% 1.7 Relevance to Course Modules
% --------------------------------

\section{Relevance to Course Modules}

The design and implementation of QVM draws directly upon core computer science coursework completed during our degree program at COMSATS University Islamabad, Vehari Campus:

\subsection{Programming Fundamentals and OOP}
Object-oriented principles govern the core codebase. Encapsulation cleanly separates the circuit definition (\texttt{QuantumCircuit}), the compilation pipeline (\texttt{Transpiler}), and the execution backends (\texttt{Simulator}, \texttt{MPSSimulator}). Polymorphism allows both simulation backends to expose identical execution interfaces, letting the caller swap between dense statevector and tensor network engines without modifying surrounding driver code.

\subsection{Data Structures and Algorithms}
Graph theory lies at the heart of quantum circuit transpilation. Physical hardware topologies are modeled as undirected graphs $G=(V, E)$. When two logical qubits need to interact across disconnected nodes, Breadth-First Search (BFS) computes the shortest path through the coupling graph. In SABRE, the router evaluates lookahead gate queues with decay weighting to select optimal candidate SWAPs. Abstract Syntax Trees (ASTs) represent parsed circuit grammars, and priority queues assist in gate dependency resolution.

\subsection{Linear Algebra and Numerical Methods}
Quantum mechanics is linear algebra in Hilbert space. Gate execution translates directly to unitary matrix-vector multiplication. Understanding Kronecker tensor products ($A \otimes B$), complex vector spaces $\mathbb{C}^{2^n}$, Hermitian operators, and Singular Value Decomposition (SVD) for tensor truncation in the MPS engine was essential for writing correct simulator routines.

\subsection{Compiler Construction and Formal Languages}
The parsing subsystem applies classic compiler concepts: lexical tokenization, Context-Free Grammars (CFGs) written in EBNF, LALR(1) shift-reduce parsing, and syntax tree traversal. The transpilation layer functions as a target-specific code generator, lowering high-level gates into hardware-supported basis gates much like a classical compiler lowers high-level expressions into machine registers.

\subsection{Software Engineering and Quality Assurance}
Throughout development, we followed iterative software engineering practices. We documented functional requirements, constructed UML behavioral and structural models, established continuous testing with 215 automated tests, and tracked defects using Git feature branches.

% --------------------------------
% 1.8 Chapter Summary
% --------------------------------

\section{Chapter Summary}

This chapter framed the central motivation behind the Quantum Virtual Machine: breaking the cycle of vendor lock-in and hardware fragmentation in modern quantum computing. By evaluating existing tools like Qiskit, Cirq, ProjectQ, and PyQuil, we identified the need for a lightweight, transparent, and educational execution platform. We detailed the primary project deliverables, explored the mathematical memory constraints inherent to classical quantum simulation, and reviewed the technology stack supporting the system. Finally, we traced the academic roots of the project back to core undergraduate computer science principles. The next chapter presents the formal problem definition, system scope, and modular architecture.
